\section{Introduction}

Consumer Android apps combine sensitive permissions, exported components, embedded SDKs, cloud services, and encrypted network traffic. Platform permission controls and Google Play Data Safety disclosures provide useful context \cite{androidPermissions2026,googlePlayDataSafety2026}, but technical claims still require evidence tied to a specific app build.

This paper investigates how build-aligned static exposure indicators and runtime network observations can jointly characterize privacy and security risk posture. Here, risk means evidence-supported opportunity for privacy or security harm, not a calibrated probability of exploitation or realized harm. The study is organized around three research questions.
\begin{itemize}
    \item RQ1: Which static exposure patterns appear in the selected app builds?
	\item RQ2: How does runtime behavior vary across strict-idle, quiescent foreground (QFG), and interactive classes for those builds?
	\item RQ3: Where do the layers agree or diverge, and how should those divergences inform review without a composite score?
\end{itemize}

A single evidence stream is insufficient. Static analysis can identify packaged exposure, but it cannot establish executed behavior. Runtime observations can capture traffic volume and interaction-driven change, but they are scenario-bounded \cite{sutter2024dynamic}. Data Safety labels and privacy-configurable SDKs are likewise not substitutes for build-specific technical evidence \cite{arkalakis2024datasafety,khandelwal2024privacylabels,zhang2024privacysdks}.

We present a 15-app study using ScytaleDroid, a database-backed platform that links APK harvests, hashes, static findings, runtime evidence packs, and quality-assurance state. The unit of analysis is a selected-build bundle rather than a mutable installed-app label. Because popular apps may update mid-wave, we use a fixed 14-day evidence window as a provenance and build-alignment control, not a longitudinal repeatability experiment. At cutoff, all 15 apps had build-backed static and runtime evidence. The main contributions are:
\begin{itemize}
    \item a 15-app static-runtime analysis with exact selected-build provenance;
    \item a recent-evidence cutoff that keeps versioned measurements valid under app-store updates;
    \item separate reporting of strict-idle, QFG, and interactive runtime classes;
    \item descriptive profiles that compare static exposure with runtime behavior as risk-posture dimensions, without an opaque composite score.
\end{itemize}
